
\section{Poset-enriched Bisupport Categories}

A poset-enriched bisupport category is a bisupport category together with a partial order $\leq$ on each hom set such that composition is monotone in each argument, such that support and cosupport are monotone and such that
each support and cosupport morphism is a subidentity. We will identify these 
additional axioms as P.1, P.2, PS.1, PC.1, PS.2 and PC.2 as follows:

\noindent [P.1]\quad
If $\coequaliser{c}{g}{\paralleldiag{a}{b}{f'}{f}}$ then 
\[
f \leq f' \;\Longrightarrow\;f\circ g \leq f'\circ g
\]
\noindent [P.2]\quad
If $\equaliser{a}{f}{\paralleldiag{b}{c}{g'}{g}}$ then 
\[
g \leq g' \;\Longrightarrow\;f\circ g \leq f\circ g'
\]

\noindent [PS.1]\quad
If $\paralleldiag{a}{b}{f}{g}$ then 
\[
f \leq g \;\Longrightarrow\;\rst{f} \leq \rst{g}.
\]
\noindent [PC.1]\quad
If $\paralleldiag{a}{b}{f}{g}$ then 
\[
f \leq g \;\Longrightarrow\;\rng{f} \leq \rng{g}.
\]

\noindent [PS.2]\quad
If $f : a \morph b$, then
\[
\rst{f} \leq id
\]
\noindent [PC.2]\quad
If $f : a \morph b$, then
\[
\rng{f} \leq id
\]

Note that strictly speaking we do not really need PC.2 as an axiom because it 
can be derived from PS.2 using xxx.

When reasoning about poset-enriched bisupport categories, we can appeal
to duality as we can within the context of bisupport categories. This is
because the poset enrichment of any category $\catcw$ induces a poset
enrichment of the opposite category $\catc^{op}$: the order on each
hom
set is unchanged, while composition is reversed.

The additional
poset-enrichment axioms are dual pairs:
$\mathrm{P.1}$ and $\mathrm{P.2}$ interchange
when passing to the opposite category; likewise,
$\mathrm{PC.1}$ is dual to
$\mathrm{PS.1}$, and $\mathrm{PC.2}$ is dual to $\mathrm{PS.2}$.

\begin{lemma}
\llabel{idempotentCompositionAndOrdering}
If \catcw is a poset-enriched bisupport category then 

(i) if $f: a \morph b$ and $g:a \morph b'$ 
then $\rst{f} \circ g \leq g$,

(ii) if $f: a \morph b$ and $g: b \morph c$ 
then  $\rng{f} \circ g \leq g$,

(iii) if $f : a \morph b$ and $g : b \morph c$ then $f \circ \rst{g} \leq f$,

(iv) if $f : a \morph b$ and $g : a' \morph b$ then $f \circ \rng{g} \leq f$,

(v) if $f : a \morph b$ and $g : b \morph c$ then $\rst{f \circ g} \leq \rst{f}$,

(vi) if $f : a \morph b$ and $g : b \morph c$ 
then $\rng{f \circ g} \leq \rng{f}$,

(vii) if $f : a \morph b$ and $f' : a \morph b'$ and $\rst{f} \leq \rst{f'}$ 
then  $\rst{f} \circ \rst{f'} = \rst{f}$,

(viii) if $f : a \morph b$ and $f' : a' \morph b$ and $\rng{f} \leq \rng{f'}$ 
then  $\rng{f} \circ \rng{f'} = \rng{g}$.

\end{lemma}
\begin{proof}
(i) Because from P.1 and PS.1 we have that
$
\rst{f} \circ g \leq id_a \circ g.
$

(ii) Follows from (i) and because $\rng{f}: b \morph b$ and
$\rst{\rng{f}}=\rng{f}$, by SC (i). 

(iii) is the dual to (ii).

(iv) is the dual to (i).

(v) \workt{???}

(vi)is the dual to (v).

(vii) $\rst{f} \circ \rst{f'} \leq  \rst{f}$ by 
lemma \lref{supportLemma} (). For the required identity it suffices to show
in addition that $\rst{f} \leq \rst{f} \circ \rst{f'}$. 
This is the case because from 
\[
\rst{f} \leq \rst{f'}
\]
it follows  by monotonicity of composition that
\[
\rst{f} \circ \rst{f} \leq \rst{f} \circ \rst{f'}
\]
which by lemma \lref{supportLemma} (iii) simplifies to
\[
\rst{f} \leq \rst{f} \circ \rst{f'}.
\]

(viii) is the dual to (vii).
\end{proof}

\begin{lemma}
\llabel{diminishingReturn}
\[
\rst{\rst{f} \circ g} \leq \rst{f}
\]
\end{lemma}
\begin{proof}
By S.3,
\[
\begin{aligned}
\rst{\rst{f}\circ g}
    &= \rst{f}\circ\rst{g}\\
    &\leq \rst{f}\circ id\\
    &= \rst{f}.
\end{aligned}
\]
The inequality follows from PS.2 and P.2.
\end{proof}
\begin{construction}
\llabel{fghConstructionRevised}
\begin{aside}
This revised version is tailored for the context of poset enriched bisupport categories. In range categories the inequality $h \leq h_L \circ h_R$ 
can be replaced by an equality. Turns out though that the inequality suffices for what we need it for working with tight zigzags. 
\end{aside}
If 
$\begin{array}{c p {0.3cm} c p{0.3cm} c }
\nudgeup{0.75cm}
          && \bOb{y} &&         \\[0.25cm]
  \bOb{x} &&         && \bOb{z} \\[0.4cm]
\end{array}
\begin{arrows}
\ncarr{x}{y}\alabel{f}
\ncarr{y}{z}\alabel{g}
\drawMorphismArcing{x}{z}{b}{h}{-15}{-15}[0.5] 
\end{arrows}
$
and $h \leq f \circ g$
in poset-enriched bisupport category \catcw
then 
if we define 
\[
h_L = \rst{h} \circ f \circ \rst{g \circ \rng{h}}
\]
and  
\[
h_R = \rng{\rst{h} \circ f} \circ g \circ \rng{h}.
\] 
then 
\begin{enumerate}[(i)]
\item
\[
h_L \leq f
\]
\item 
\[
h_R \leq g
\]
\item
\[
\rst{h_L}=\rst{h}
\]
\item
\[
\rng{h_R}=\rng{h}.
\]
\item 
\[
h \leq h_L \circ h_R
\]
\end{enumerate}
\end{construction}
\begin{proof}\leavevmode\par
\begin{enumerate}[(i)]
\item %(i)
That $h_L \leq f$
is the case
follows from the form the definition
of $h_L$  as restriction idempotents
composed with $f$,
see lemma \lref{idempotentCompositionAndOrdering}.

\item %(ii)
Follows by the same argument as (i).

\item %(iii)
To prove that $\rst{h_L}=\rst{h}$ 
we prove first that $\rst{h_L}\leq \rst{h}$
and second that $\rst{h} \leq \rst{h_L}$.

For the first we have 
\begin{align*}
\rst{h_L}
    &= \rst{\rst{h} \circ f \circ \rst{g \circ \rng{h}}}
         && \mbox{by definition of $h_L$,}\\
    &= \rst{\rst{h} \circ f \circ g \circ \rng{h}}
        && \mbox{by S.4,}\\
    &\leq \rst{h} 
        && \mbox{by lemma \lref{diminishingReturn}.}
\end{align*}

For the second we have
\begin{align*}
\rst{h} &= \rst{\rst{h} \circ h \circ \rng{h}}
              &&\mbox{by S.1 and C.1,} \\
        &\leq \rst{\rst{h} \circ f \circ g \circ \rng{h}}
            && \mbox{because $h \leq f \circ g$ and by monotonicity,} \\
        &= \rst{\rst{h} \circ f \circ \rst{g \circ \rng{h}}}
            && \mbox{S.4} \\
        &= \rst{h_L} && \mbox{by definition.}
\end{align*}
\item %(iv) 
The identity $\rng{h_R}=\rng{h}$ 
follows by duality because $h_R$, as defined, is the dual of $h_L$ 
and because the dual of $\rst{h}$ is $\rng{h}$.
\item %(v)
The inequality $h \leq h_L \circ h_R$ follows 
because 
\begin{align*}
h &= \rst{h} \circ h \circ \rng{h} 
               && \mbox{by S.1 and C.1,} \\
  &\leq \rst{h} \circ f \circ g \circ \rng{h}
                && \mbox{by monotonicity of composition and since $h \leq f \circ g$,} \\
  &= \rst{h} \circ f \circ \rng{\rst{h} \circ f} 
       \circ 
     \rst{g \circ \rng{h}} \circ g \circ \rng{h}
                && \mbox{by S.1 and C.1,} \\
  &= \rst{h} \circ f \circ \rst{g \circ \rng{h}} 
       \circ
     \rng{\rst{h} \circ f} \circ g \circ \rng{h}
                 && \mbox{\parbox[t]{5cm}{because support idempotents
                             commute (S.2) 
                             and because $\rng{\rst{h} \circ f}$ 
                             is a support idempotent (SC(i)), }}\\
  &= h_L \circ h_R && \mbox{from how $h_L$ and $h_R$ are defined.}
\end{align*}
\end{enumerate}
\end{proof}

